\scanpage{39}
\mspage{- 43 -}
n\textsuperscript{o} 272

dans $A'$ des éléments $y_\alpha$ ($\alpha=1,\ldots,s$) dont les images $\bar y_\alpha$ dans $L_{A'}$ sont linéairement indépendantes sur $K_A$ ; on a évidemment $v'(y_\alpha)=0$ pour tout $\alpha$. Nous allons montrer que les éléments $x_i y_\alpha$ sont linéairement indépendants sur $K$, ce qui établira bien l'inégalité $rs\leq n$.

Supposons donc qu'il existe une relation linéaire non triviale de la forme ;
\begingroup\csname tagsleft@true\endcsname
\[
\sum a_{i\alpha}x_i y_\alpha=0\quad,\ a_{i\alpha}\in K. \tag*{(2)}
\]
\endgroup

Soit $a_{j\beta}x_j y_\beta$ celui des termes de (2) dont la valuation est minimum ; on a donc $v'(a_{i\alpha}x_i y_\alpha)\geq v'(a_{j\beta}x_j y_\beta)$ pour tout couple $(i,\alpha)$ ; si $i\ne j$, on ne peut avoir égalité, car cela entraînerait
\begingroup\csname tagsleft@true\endcsname
\[
v'(x_i)-v'(x_j)=v'(a_j)-v'(a_i)\qquad\equiv 0\ \mathrm{mod.}\ \Gamma_v, \tag*{(3)}
\]
\endgroup
contrairement au choix des $x_i$. On a donc :
\[
v'(a_{i\alpha}x_i y_\alpha)>v'(a_{j\beta}x_j y_\beta)\qquad\text{pour } i\ne j
\]

En multipliant la relation (2) par $(a_{j\beta}x_j)^{-1}$ et en posant $b_\alpha=a_{j\alpha}\cdot a_{j\beta}^{-1}$ on obtient alors une relation de la forme :
\begingroup\csname tagsleft@true\endcsname
\[
\sum b_\alpha y_\alpha+z=0\quad,\ \text{avec }v'(b_\alpha)\geq 0,\ v^{\unsure{\prime|:}}(z)>0, \tag*{(4)}
\]
\endgroup
d'où, dans $L_{A'}$, une relation :
\begingroup\csname tagsleft@true\endcsname
\[
\sum\bar b_\alpha\bar y_\alpha=0. \tag*{(5)}
\]
\endgroup

Comme $b_\beta=1$, ceci contredit l'hypothèse faite sur les $y_\alpha$, et achève la démonstration.

PROPOSITION 2 - \uline{Supposons que} $L$ \uline{soit une extension algébrique de} $K$. \uline{Alors le groupe quotient} $\Gamma_{v'}/\Gamma_v$ \uline{est un groupe de torsion}, \uline{et l'extension} $L_{A'}/K_A$ \uline{est une extension algébrique}.

Soit $\{L_\omega\}$ la famille des extensions finies de $K$ contenues dans $L$, et soit $\Gamma_\omega=v'(L_\omega^*)$. Le groupe $\Gamma_{v'}$ est réunion de la famille filtrante croissante formée par les $\Gamma_\omega$ ; comme, d'après la prop.1; les groupes
% archive sheet number: 38
